(1) 反射镜平移 $\Delta x$ 时光往返多走 $2\Delta x$，故光程差变化 $\delta\Delta=2\Delta x$。单色光相位变化 $\delta\phi=2\pi\delta\Delta/\lambda_0=4\pi\Delta x/\lambda_0$，对应中心频率 $\lambda_0=c/\nu_0$。这里已把两次通过一段光程计入，不能只写单程位移。
(2) 固定探测点连续两个亮纹条件的相位差为 $2\pi$，所以 $4\pi\Delta x/\lambda_0=2\pi$，得到每条亮纹对应镜位移 $\Delta x=\lambda_0/2$。若镜向相反方向移动，条纹移动方向反转但绝对通过数仍由路径改变量决定。
(3) 按 Wiener–Khinchin 相干关系，$\gamma(\Delta)=\int S(\nu)e^{i2\pi\nu\Delta/c}d\nu/\int S(\nu)d\nu$。高斯傅里叶积分给 $\gamma(\Delta)=\exp[i2\pi\nu_0\Delta/c]\exp[-2\pi^2\sigma_\nu^2\Delta^2/c^2]$，其中第一指数是快条纹相位，第二指数是相干包络衰减。
(4) 两个臂强度分别 $I_1,I_2$ 的干涉可见度为 $\mathcal V=2\sqrt{I_1I_2}|\gamma|/(I_1+I_2)$。在等强条件下它化简为 $\mathcal V(\Delta)=\exp[-2\pi^2\sigma_\nu^2\Delta^2/c^2]$；令可见度为 $1/e$ 得 $\ell_c=c/(\sqrt2\pi\sigma_\nu)$。若两臂不等强，即使 $\Delta=0$ 也不能达到单位可见度。
